\begin{resumo}[Abstract]
\begin{SingleSpace}
Investors use market indices as a benchmark for return analysis, but these indices are generalist, reflecting average behavior. Building indices requires the collection, processing, and composition of metrics over a database. This work proposes a dimensional model representing data from stocks listed on the US stock market, and an environment to enable the creation of new indices that more adequately reflect different investor profiles. The evaluation was conducted by comparing composed portfolios against the S\&P 500 index, demonstrating a performance gain when modifiers are incorporated into the selection criteria. The main contribution of this study is the integration of data collection, modeling, and metric composition, rather than merely generating an isolated indicator. The curation of the asset set and the materialization of derived metrics in the data layer ensure the reproducibility of the results and the auditability of the analytical workflow.
\end{SingleSpace}
\vspace{\onelineskip}
\textbf{Keywords}: risk analysis; stock market; data mart; dimensional modeling; smart beta.
\end{resumo}